% GENERATED FILE. Edit canonical lessons, then run npm run generate:books.
% canonical-source-hash: fnv1a64:f15cc20efcc1ae2a
% canonical-lessons: GU-C157-gothavvu, GU-C157-mukvu, GU-C157-tagvu, GU-C157-paddo, GU-C157-chajli

\chapter{Putting things in their place}
\label{ch:gu-putting-things-in-their-place}
\hlchaptermodality{157}

\begin{chapteropening}
\textbf{By the end of this chapter:} \emph{I can say where I put and hang things at home.}

Complete the last lesson of chapter 157: I can say where I put and hang things at home.
\end{chapteropening}

\section[\texorpdfstring{\gu{ગોઠવવું} (goṭhavvũ)}{ગોઠવવું (goṭhavvũ)}]{\gu{ગોઠવવું} (goṭhavvũ) — to arrange, to set out}

\label{lesson:GU-C157-gothavvu}



\begin{quote}
Before the new one: say the Gujarati for to decorate, then the Gujarati for to tie.
\end{quote}

\subsection*{You'll want to know: \gu{ગોઠવવું}}
\textbf{\gu{ગોઠવવું}} — \emph{goṭhavvũ} — \textquotedblleft{}to arrange, to set out\textquotedblright{}.

\begin{quote}
\textbf{\gu{તે} \gu{પુસ્તકો} \gu{ગોઠવે} \gu{છે}.} — \emph{te pustako goṭhave che} — \textquotedblleft{}He arranges the books.\textquotedblright{}
\end{quote}

\begin{cousinweb}
The origin of \textbf{\gu{ગોઠવવું}} (\emph{goṭhavvũ}) is not settled. CDIAL lists no Gujarati form of it, and EtymDB has no entry for it.
\end{cousinweb}

\begin{grammarlens}[title={What you've built}]
One more word for everyday life at A2.
\end{grammarlens}

\subsection*{Your turn}
\begin{itemize}
\raggedright
  \item \emph{Say it:} \emph{goṭhavvũ}
  \item \emph{Say it:} \emph{goṭhavvũ}, once more
  \item \emph{Say it:} \emph{te pustako goṭhave che} — He arranges the books.
  \item \emph{Recall:} say the Gujarati for to decorate, then the Gujarati for to tie, then say \emph{goṭhavvũ} again
\end{itemize}

\begin{tcolorbox}[breakable,colback=teal!4,colframe=teal!35!black,title={Before you move on}]
What does \gu{ગોઠવવું} mean? (\textquotedblleft{}To arrange, to set out\textquotedblright{}.) Say it once more.
\end{tcolorbox}
\section[\texorpdfstring{\gu{મૂકવું} (mūkvũ)}{મૂકવું (mūkvũ)}]{\gu{મૂકવું} (mūkvũ) — to put, to place}

\label{lesson:GU-C157-mukvu}



\begin{quote}
Before the new one: say the Gujarati for to tie, then the Gujarati for to arrange.
\end{quote}

\subsection*{You'll want to know: \gu{મૂકવું}}
\textbf{\gu{મૂકવું}} — \emph{mūkvũ} — \textquotedblleft{}to put, to place\textquotedblright{}.

\begin{quote}
\textbf{\gu{હું} \gu{ચાવી} \gu{કબાટમાં} \gu{મૂકું} \gu{છું}.} — \emph{hũ cāvī kabāṭmā̃ mūkũ chũ} — \textquotedblleft{}I put the key in the cupboard.\textquotedblright{}
\end{quote}

\begin{cousinweb}
\textbf{\gu{મૂકવું}} (\emph{mūkvũ}) is a tadbhava. It comes from \emph{*mukna}, \textquotedblleft{}loosed\textquotedblright{}, a reconstructed form beside Sanskrit \emph{mukta}, \textquotedblleft{}released\textquotedblright{}; CDIAL glosses the Gujarati verb \textquotedblleft{}to let go, put, place\textquotedblright{} (CDIAL 10157).
\end{cousinweb}

\begin{grammarlens}[title={What you've built}]
One more word for everyday life at A2.
\end{grammarlens}

\subsection*{Your turn}
\begin{itemize}
\raggedright
  \item \emph{Say it:} \emph{mūkvũ}
  \item \emph{Say it:} \emph{mūkvũ}, once more
  \item \emph{Say it:} \emph{hũ cāvī kabāṭmā̃ mūkũ chũ} — I put the key in the cupboard.
  \item \emph{Recall:} say the Gujarati for to tie, then the Gujarati for to arrange, then say \emph{mūkvũ} again
\end{itemize}

\begin{tcolorbox}[breakable,colback=teal!4,colframe=teal!35!black,title={Before you move on}]
What does \gu{મૂકવું} mean? (\textquotedblleft{}To put, to place\textquotedblright{}.) Say it once more.
\end{tcolorbox}
\section[\texorpdfstring{\gu{ટાંગવું} (ṭā̃gvũ)}{ટાંગવું (ṭā̃gvũ)}]{\gu{ટાંગવું} (ṭā̃gvũ) — to hang up}

\label{lesson:GU-C157-tagvu}



\begin{quote}
Before the new one: say the Gujarati for to arrange, then the Gujarati for to put.
\end{quote}

\subsection*{You'll want to know: \gu{ટાંગવું}}
\textbf{\gu{ટાંગવું}} — \emph{ṭā̃gvũ} — \textquotedblleft{}to hang up\textquotedblright{}.

\begin{quote}
\textbf{\gu{તે} \gu{કોટ} \gu{કબાટમાં} \gu{ટાંગે} \gu{છે}.} — \emph{te koṭ kabāṭmā̃ ṭā̃ge che} — \textquotedblleft{}He hangs his coat in the cupboard.\textquotedblright{}
\end{quote}

\begin{cousinweb}
\textbf{\gu{ટાંગવું}} (\emph{ṭā̃gvũ}) is a tadbhava. It comes from \emph{*ṭaṅg}, \textquotedblleft{}hang up\textquotedblright{}, a reconstructed form, and CDIAL glosses the Gujarati verb the same way (CDIAL 5436).
\end{cousinweb}

\begin{grammarlens}[title={What you've built}]
One more word for everyday life at A2.
\end{grammarlens}

\subsection*{Your turn}
\begin{itemize}
\raggedright
  \item \emph{Say it:} \emph{ṭā̃gvũ}
  \item \emph{Say it:} \emph{ṭā̃gvũ}, once more
  \item \emph{Say it:} \emph{te koṭ kabāṭmā̃ ṭā̃ge che} — He hangs his coat in the cupboard.
  \item \emph{Recall:} say the Gujarati for to arrange, then the Gujarati for to put, then say \emph{ṭā̃gvũ} again
\end{itemize}

\begin{tcolorbox}[breakable,colback=teal!4,colframe=teal!35!black,title={Before you move on}]
What does \gu{ટાંગવું} mean? (\textquotedblleft{}To hang up\textquotedblright{}.) Say it once more.
\end{tcolorbox}
\section[\texorpdfstring{\gu{પડદો} (paḍdo)}{પડદો (paḍdo)}]{\gu{પડદો} (paḍdo) — curtain}

\label{lesson:GU-C157-paddo}



\begin{quote}
Before the new one: say the Gujarati for to put, then the Gujarati for to hang up.
\end{quote}

\subsection*{You'll want to know: \gu{પડદો}}
\textbf{\gu{પડદો}} — \emph{paḍdo} — \textquotedblleft{}curtain\textquotedblright{}.

\begin{quote}
\textbf{\gu{બારીનો} \gu{પડદો} \gu{ભૂરો} \gu{છે}.} — \emph{bārīno paḍdo bhūro che} — \textquotedblleft{}The window curtain is blue.\textquotedblright{}
\end{quote}

\begin{cousinweb}
\textbf{\gu{પડદો}} (\emph{paḍdo}) is a loan, usually traced to Persian \emph{parda}, \textquotedblleft{}curtain, veil\textquotedblright{}. CDIAL, which follows inherited words, lists no Gujarati form.
\end{cousinweb}

\begin{grammarlens}[title={What you've built}]
One more word for everyday life at A2.
\end{grammarlens}

\subsection*{Your turn}
\begin{itemize}
\raggedright
  \item \emph{Say it:} \emph{paḍdo}
  \item \emph{Say it:} \emph{paḍdo}, once more
  \item \emph{Say it:} \emph{bārīno paḍdo bhūro che} — The window curtain is blue.
  \item \emph{Recall:} say the Gujarati for to put, then the Gujarati for to hang up, then say \emph{paḍdo} again
\end{itemize}

\begin{tcolorbox}[breakable,colback=teal!4,colframe=teal!35!black,title={Before you move on}]
What does \gu{પડદો} mean? (\textquotedblleft{}Curtain\textquotedblright{}.) Say it once more.
\end{tcolorbox}
\section[\texorpdfstring{\gu{છાજલી} (chājlī)}{છાજલી (chājlī)}]{\gu{છાજલી} (chājlī) — shelf}

\label{lesson:GU-C157-chajli}



\begin{quote}
Before the new one: say the Gujarati for to hang up, then the Gujarati for curtain.
\end{quote}

\subsection*{You'll want to know: \gu{છાજલી}}
\textbf{\gu{છાજલી}} — \emph{chājlī} — \textquotedblleft{}shelf\textquotedblright{}.

\begin{quote}
\textbf{\gu{છાજલીમાં} \gu{ત્રણ} \gu{પુસ્તકો} \gu{છે}.} — \emph{chājlīmā̃ traṇ pustako che} — \textquotedblleft{}There are three books on the shelf.\textquotedblright{}
\end{quote}

\begin{cousinweb}
The origin of \textbf{\gu{છાજલી}} (\emph{chājlī}) is not settled. CDIAL lists no Gujarati form of it; it lists Gujarati \emph{chāj}, \textquotedblleft{}thatch\textquotedblright{}, and \emph{chājũ}, \textquotedblleft{}thatched lean-to shed\textquotedblright{}, but does not link them to a shelf (CDIAL 5023).
\end{cousinweb}

\begin{grammarlens}[title={What you've built}]
One more word for everyday life at A2.
\end{grammarlens}

\subsection*{Your turn}
\begin{itemize}
\raggedright
  \item \emph{Say it:} \emph{chājlī}
  \item \emph{Say it:} \emph{chājlī}, once more
  \item \emph{Say it:} \emph{chājlīmā̃ traṇ pustako che} — There are three books on the shelf.
  \item \emph{Recall:} say the Gujarati for to hang up, then the Gujarati for curtain, then say \emph{chājlī} again
\end{itemize}

\begin{tcolorbox}[breakable,colback=teal!4,colframe=teal!35!black,title={Before you move on}]
What does \gu{છાજલી} mean? (\textquotedblleft{}Shelf\textquotedblright{}.) Say it once more.
\end{tcolorbox}
